\documentclass[conference]{IEEEtran}
\IEEEoverridecommandlockouts

\usepackage{cite}
\usepackage{amsmath,amssymb}
\usepackage{graphicx}
\usepackage{booktabs}
\usepackage{url}
\usepackage[hidelinks]{hyperref}

\begin{document}

\title{Batch Tiling Overstates the Cost of Causal\\Sleep Staging: Evidence from Two Datasets}

\author{\IEEEauthorblockN{Jahnvi R, Swapnil S, Shreeya S Methuku, Shreshtha Dixit}
\IEEEauthorblockA{Department of Computer Science (AIML) and Engineering\\
PES University, Bengaluru, India\\
\texttt{jahnvir04@gmail.com}, \texttt{swapnil160305@gmail.com},\\
\texttt{ft.shreeyyymethuku@gmail.com}, \texttt{shreshtha.dixit.511@gmail.com}}}

\maketitle

\begin{abstract}
A sleep-staging model that must run on a wearable device may use only past signal, and the
usual assumption is that this constraint costs accuracy. We measure that cost with a matched
ablation in which the same network is trained with and without access to future signal,
holding architecture, parameter count, data, folds and seeds fixed, and we run the identical
experiment on two datasets: 78 Sleep-EDF subjects and the 25 subjects of DOD-H. Measured the
way this literature normally evaluates, by tiling a recording with non-overlapping context
windows, the causal constraint costs $0.025$ and $0.033$ Cohen's $\kappa$ respectively, and
on both datasets it appears to fall on REM. We show that most of that number is produced by
the evaluation protocol rather than by the constraint. Tiling scores predictions at every
position of the context buffer, including positions with almost no history, which a streaming
system passes through once at start-up and never again; the causal model is crippled there
and the non-causal model is not, because it still has lookahead. Switching to a streaming
protocol therefore lifts the two arms by unequal amounts, and the excess lift given to the
causal arm is $+0.0314\,\kappa$ on Sleep-EDF-78 and $+0.0292$ on DOD-H, positive in 30 of 30
paired folds, with corrected $t(14) = 12.69$ and $6.48$. That asymmetry is the paper's result,
and it is what replicates. Re-measured under the deployment protocol the causality cost is
$+0.0065$ and $-0.0037\,\kappa$, neither resolvable, and the apparent REM concentration
falls from $0.060$ F1 to $0.002$ and $0.023$. We also identify and close a look-ahead leak in
the per-epoch normalization used here and in most published baselines; the causal replacement
improves both arms, so it was never a trade. The model is strictly causal by construction,
verified bit-for-bit from raw sample to logit, emits one label per second, and runs in
30{,}757 parameters at $0.026$~ms per second of EEG on a consumer CPU. Reported causality
penalties in this literature should be read with the evaluation protocol in mind.
\end{abstract}

\begin{IEEEkeywords}
sleep staging, electroencephalography, causal models, temporal convolutional networks,
self-attention, real-time inference, evaluation methodology
\end{IEEEkeywords}

\section{Introduction}

Standard clinical polysomnography (PSG) is scored in 30-second epochs by a human technician
following American Academy of Sleep Medicine (AASM) conventions~\cite{aasm}, and most
automated sleep-staging models mirror that unit: they classify one 30-second window at a time
and are often free to use signal from later in the recording, or the entire night, to label an
earlier epoch. This is a reasonable design for offline analysis of a completed overnight study.
It is unsuitable for a wearable device that must report a stage as it happens: for real-time
alerting, closed-loop stimulation, or a live dashboard, a model cannot wait for a later epoch
before deciding what occurred a minute earlier.

Causality is a real constraint, so the natural question is what it costs. That question is
rarely answered cleanly. Published causal and non-causal staging models differ in capacity,
data, preprocessing and optimization all at once, and the gap between them is not attributable
to causality in particular. We answer it directly with a matched ablation: the same
architecture, parameter count, data, folds and seeds, with only the left-padding and the
attention mask removed.

The first result of that experiment is the one we expected. Under the evaluation protocol this
literature normally uses, the causal constraint costs $0.025\,\kappa$ on Sleep-EDF-78 and
$0.033$ on DOD-H, consistently across seeds. The second result is that most of that number is
not a property of the constraint. It is produced by scoring the model at context-buffer
positions that a deployed system occupies once, at start-up, and understanding why turns out
to be the more useful contribution.

Our contributions are three, and the second and third are consequences of the first rather
than independent claims.

\emph{First, the evaluation protocol.} We show that batch-tiled evaluation systematically
penalizes causal models, and we measure the penalty as a within-arm quantity that does not
depend on the two arms being comparable. Moving the same trained weights from the tiled to the
streaming protocol raises the causal arm by $0.0350\,\kappa$ on Sleep-EDF-78 and the non-causal
arm by only $0.0036$; on DOD-H the same figures are $0.0345$ and $0.0052$. The excess lift
given to the causal arm is $+0.0314$ and $+0.0292\,\kappa$, positive in all 30 paired folds
across the two datasets, with Nadeau--Bengio corrected $t(14) = 12.69$ ($p = 4.6\times10^{-9}$)
and $t(14) = 6.48$ ($p = 1.4\times10^{-5}$). We verify the mechanism arithmetically: the tiled
score is the mean of a curve over buffer positions, reproduced to within $0.003\,\kappa$ for
both arms on both datasets, and the penalty is generated almost entirely by positions carrying
under 30 seconds of history.

\emph{Second, what causality actually costs once that artifact is removed.} Under the
deployment protocol the measured cost is $+0.0065\,\kappa$ on Sleep-EDF-78 and $-0.0037$ on
DOD-H, neither of which is resolvable, and the apparent concentration of the penalty in REM,
$0.060$ F1 under tiling on both datasets, falls to $0.002$ and $0.023$. We give the full
latency--accuracy curve, which states what lookahead is worth in seconds of permitted delay,
and we report the one place where the two datasets disagree rather than the one place they
agree: within a fixed 120-second buffer the non-causal arm never exceeds the causal arm's
zero-latency score on Sleep-EDF, but does exceed it by $0.014\,\kappa$ at 46 seconds of delay
on DOD-H.

\emph{Third, causality as a property of the pipeline.} Per-epoch $z$-scoring, close to
universal in this literature, normalizes an early sample using statistics computed over the
whole epoch, including its future. Replacing it with a trailing-window statistic computable one
sample at a time makes the path from raw sample to logit end-to-end causal, verified
bit-for-bit, and improves both arms of the ablation, by $0.024\,\kappa$ causal and $0.021$
non-causal. A pipeline step that reads the future is not buying accuracy in exchange for
causality here. Alongside this we report the model itself: a strictly causal staging network
whose causality is verified by perturbation test rather than asserted by construction, emitting
one label per second in 30{,}757 parameters, with a parameter-free causal smoothing step that
adds $0.010\,\kappa$ at inference with no retraining.

\section{Related Work}

Automated sleep staging from single- or few-channel EEG has moved from hand-engineered spectral
features to end-to-end deep learning over the past decade. Supratak et al.\ proposed
DeepSleepNet~\cite{deepsleepnet}, which pairs a dual-branch CNN with a bidirectional LSTM
modelling transition rules across epochs. It reports strong accuracy on Sleep-EDF but reads both
past and future epochs through its bidirectional recurrence and uses roughly 22 million
parameters. AttnSleep~\cite{attnsleep} replaces the recurrent stage with multi-head attention
over a sequence of epochs, cutting parameter count by about three orders of magnitude while
remaining competitive; its attention is not masked, and its 30-second epoch length makes it a
batch rather than a streaming method. CareSleepNet~\cite{caresleepnet} is a more recent hybrid
convolution--transformer that additionally fuses EEG and EOG, at roughly 22.7 million
parameters, and is likewise not causal. Other established single-channel architectures include
SeqSleepNet~\cite{seqsleepnet}, TinySleepNet~\cite{tinysleepnet}, the single-epoch CNN of Sors
et al.~\cite{sors}, SleepEEGNet~\cite{sleepeegnet} and IITNet~\cite{iitnet}. All of these read
future signal, label at 30-second granularity, or both.

\subsection{High-Frequency Staging Is Already Solved, Offline}

U-Sleep~\cite{usleep} is the closest prior work on output resolution, and the clearest statement
of why resolution by itself is not the contribution here. Trained and evaluated on 15{,}660
participants across 16 clinical studies, it reports classification performance at fourteen output
frequencies spanning 2 to 7{,}680 stages per minute, the latter being one label per sample at
128~Hz. Its predecessor U-Time~\cite{utime} introduced the mechanism: an encoder--decoder
classifies every input sample, and a small segment classifier aggregates those classifications
over whatever interval is asked for, so the training label rate and the inference output rate are
independent. High-frequency sleep staging is therefore not an open problem, and it has not been
one since 2019.

It is, however, solved non-causally. U-Sleep is a bidirectional encoder--decoder whose receptive
field spans up to 13.5 minutes, roughly half of which lies after the second being labelled, and it
makes no streaming or real-time claim. The question this paper asks is what the same resolution
costs when the model may not look forward at all.

U-Sleep also supplies the motivation for asking. Generating stages at those fourteen frequencies
and testing how well the resulting hypnograms separate obstructive sleep apnea patients from
healthy controls, the authors report mean F1 rising from 0.60 at 2 stages per minute to 0.94 at
1{,}280, and conclude that high-frequency stages are clinically informative. Fine-grained staging
is worth having. Whether it can be had causally, and at what price, is what we measure.

\subsection{Causal and Real-Time Staging}

A separate line of work targets causal or real-time operation directly. Bresch et
al.~\cite{bresch} apply recurrent deep networks to real-time single-channel staging. Pan et al.'s
CausalAttenNet~\cite{causalattennet} is a fast causal long-term-temporal attention network for
single-channel EEG, and reports N1 as its weakest class, consistent with the pattern observed
here. Mehrotra et al.'s SnoozeNet~\cite{snoozenet} applies causality-driven transformers to
interpretable staging, and Liu et al.'s Micro SleepNet~\cite{microsleepnet} targets real-time
staging on mobile and wearable hardware. These establish that network-level causal staging is an
active and viable direction, and this work does not claim priority over them on that point.

Streaming causal inference over raw EEG is also active beyond sleep staging.
CaMBRAIN~\cite{cambrain} applies a causal state-space model with a persistent hidden state to
artifact, abnormality, seizure and stress detection at 16 predictions per second, reporting large
efficiency gains over windowed baselines. It does not address sleep staging and does not compare
against a matched bidirectional control, so it establishes the direction rather than the
measurement attempted here.

To our knowledge none of the above reports a matched ablation in which the identical architecture
is trained with and without access to future signal on identical folds and seeds, which is what
licenses attributing a performance gap to causality rather than to capacity or data. Bresch et
al.~\cite{bresch} adopt causal recurrence by design and report no non-causal counterpart; Micro
SleepNet~\cite{microsleepnet} ablates its attention and dilation components but not causality, and
declines sequence-level modelling outright on memory grounds. We have not been able to inspect the
methods sections of~\cite{causalattennet} and~\cite{snoozenet} directly, so we state this as a
limit of our reading of this literature rather than as an exhaustive claim.

\subsection{Two Points That Cut Across These Groups}

The first concerns preprocessing. Causality is a property of a pipeline, not only of a network,
and per-epoch normalization, which computes an epoch's mean and standard deviation over all 30
seconds and applies them to each sample within it, silently reads the future. It is close to
universal in this literature. Section~\ref{sec:norm} states what it does and what closing it
costs; the answer is that closing it is free.

The second concerns evaluation. Tiling a recording with non-overlapping context windows, the
default in every implementation we have examined including our own earlier work, scores a model at
buffer positions a streaming system does not occupy after start-up. Sections \ref{sec:protocol}
and \ref{sec:latency} show that this alone accounts for most of the causality penalty we and, by
implication, others would otherwise report, on both datasets studied here.

Temporal convolutional networks with causal, dilated convolutions have been shown to match
recurrent networks on long sequence-modelling tasks while remaining strictly
forward-looking~\cite{tcn}; that building block is adopted here and combined with causally masked
attention~\cite{attention}, giving long context without the loss of parallelism that recurrence
introduces.

\section{Data and Preprocessing}

\subsection{Datasets}
\label{sec:datasets}

Two datasets are used, and the second exists to test whether the protocol finding is a property of
the first.

\textbf{Sleep-EDF Expanded} (sleep-cassette subset), distributed by PhysioNet~\cite{physionet} and
described by Kemp et al.~\cite{kemp}: a single EEG channel (Fpz-Cz) resampled to 100~Hz, five
AASM-style classes~\cite{aasm} (Wake, N1, N2, N3, REM; N3 and N4 merged per convention). Two
subsets are reported: Sleep-EDF-20 (20 subjects, 39 recordings) and Sleep-EDF-78 (78 subjects, 153
recordings, 1{,}629 hours after Wake trimming), which spans ages 25--101 and is the harder, more
representative population.

\textbf{DOD-H}~\cite{dod}, the healthy-volunteer arm of the Dreem Open Datasets: 25 subjects,
205.5 hours, recorded at 250~Hz and scored by five independent technicians, of whom we use the
released consensus hypnogram. Its class balance differs substantially from wake-trimmed Sleep-EDF
(Wake 12.3\%, N1 6.1\%, N2 48.2\%, N3 14.2\%, REM 19.2\%), which is part of why it is a useful
second test.

All splits on both datasets use subject-wise 5-fold cross-validation, verified by unit test
(Table~\ref{tab:verified}) so that no subject's data appears in both a training and a held-out
fold.

\subsection{Porting the Pipeline to DOD-H}
\label{sec:dodh-prep}

Nothing about the architecture, the training recipe or the two evaluation protocols changes
between datasets; only the source of the EEG does. Three dataset-level decisions are forced by
DOD-H and each is recorded in the released preprocessing manifest.

\emph{Channel.} DOD-H carries twelve derivations. We use F3-M2, which is one of only four frontal
derivations present in both DOD-H and its clinical companion DOD-O, so that the same choice
extends to DOD-O without a second decision. This is not Fpz-Cz, and no cross-dataset transfer is
attempted or implied: each dataset is trained and evaluated entirely within itself.

\emph{Sampling rate.} The MRCNN kernels are defined in samples (50 and 400, that is 0.5~s and 4~s
at 100~Hz). We resample DOD-H from 250 to 100~Hz with a polyphase anti-aliasing filter, an exact
$2/5$ ratio, rather than retuning the kernels, so the network is byte-for-byte the same in both
experiments.

\emph{Wake trimming.} Sleep-EDF cassette recordings run about 20 hours with Wake at 68\%, which is
why that pipeline trims. DOD-H records are 5--10 hours with Wake already at 12.3\%, so no trimming
is applied; applying the Sleep-EDF setting would delete real sleep-adjacent Wake rather than
daytime recording.

Two further points are worth stating because they are the kind of thing that silently invalidates
a second-dataset result. Amplitude units differ between the two corpora, and the trailing $z$-score
of Section~\ref{sec:norm} divides by a trailing standard deviation, so the pipeline is
scale-invariant; we check rather than assume this, and the normalized standard deviation across the
25 DOD-H subjects lies in $[1.02, 1.12]$. And normalization is applied to the whole continuous
recording \emph{before} unscored epochs are dropped, exactly as on Sleep-EDF, because normalizing
each surviving stretch separately would reset the trailing statistic at every gap, which a device
would not do. Windows are additionally refused if they would span a dropped gap or a night join.
The focal-loss class weights are recomputed per training fold, so DOD-H's different class balance
is handled automatically rather than by reusing Sleep-EDF's weights.

\subsection{Wake Trimming on Sleep-EDF}

Sleep-EDF cassette recordings run approximately 20 hours per recording, and most subjects
contribute two consecutive night recordings, so a subject supplies roughly 40 hours of which only a
fraction is sleep. Most of the remainder is awake time outside the sleep period, and Wake
constitutes 68\% of the raw data if retained in full. Each night is trimmed to 30 minutes of Wake on
either side of that night's sleep period, applied per recording before a subject's two recordings
are concatenated, which also removes the daytime gap between them. Trimming reduces the 153
recordings of Sleep-EDF-78 from roughly 3{,}000 hours to the 1{,}629 hours reported throughout this
paper; every hour figure quoted here is post-trim. This reduces Wake to approximately 15\% of the
data, matches the convention used by the baselines in Section~\ref{sec:baselines}, and materially
changes which stages the model is rewarded for classifying correctly
(Section~\ref{sec:preproc-ablation}).

\subsection{Label Supervision}

Expert labels are per-30-second epoch on both datasets, as is standard for AASM
scoring~\cite{aasm}; no public dataset currently provides genuinely per-second expert annotation.
Each epoch's label is replicated across its 30 constituent one-second training targets, so the
network is supervised at 1~Hz even though the ground truth is coarser. Per-second accuracy should
therefore be read as an approximation. $\kappa$, and accuracy re-aggregated to 30~s by majority
vote, are the more comparable figures against 30-second baselines, and both are reported throughout.

\sloppy
\subsection{Normalization: Closing a Look-Ahead Leak}
\label{sec:norm}

Earlier versions of this pipeline $z$-scored each 30-second epoch using that epoch's own mean and
standard deviation. The network is causal with respect to its input tensor, but the pipeline that
produces that tensor is not: a normalized sample near the start of an epoch depends on raw samples
up to 30~s later in the same epoch. The leak is bounded, since it never crosses an epoch boundary,
but it is real, and it is inherited from the convention used by the baselines this work compares
against.

Closing it requires replacing the per-epoch statistic with a trailing statistic computable one
sample at a time, in which the statistics at sample $t$ are drawn only from a preceding window of
$w$ seconds and reset at the start of every recording:
\begin{equation}
\mu_t = \frac{1}{w}\sum_{i=t-w+1}^{t} x_i
\end{equation}
\begin{equation}
z_t = \frac{x_t - \mu_t}{\sigma_t + \epsilon}
\end{equation}
where $x_t$ is the raw sample at time $t$, $w$ is the window in samples, $\sigma_t$ is the trailing
standard deviation over the same window, and $\epsilon$ is a small constant guarding against
division by zero. The statistic resets at the start of each recording, so the opening seconds run
on a partial window: before a full window has been observed ($t < w$) the statistics are taken over
the samples seen so far. That is deliberate. A device powering on has no history, and the behaviour
is bounded, since within a window of $n$ samples the largest attainable standardized magnitude is
$\sqrt{n-1}$, so neither an infinity nor a NaN can arise there. This is mathematically a running
statistic a device could compute online, and \texttt{src/data/normalization.py} ships both the
offline array implementation and the sample-at-a-time filter; they agree to $5.2\times10^{-14}$
over 24{,}000 samples, so the equivalence is checked rather than asserted.

Every result reported in this paper, on both datasets, is produced under (1)--(2). The archived
per-epoch runs are retained as a controlled comparison rather than as the headline, because both
pipelines were trained from scratch with identical subjects, folds, seeds, architecture and
hyperparameters, so the difference between them isolates the normalizer.

Closing the leak improves the model. Cohen's $\kappa$ rises from $0.6406$ to $0.6649$ for the
causal arm and from $0.6693$ to $0.6898$ for the non-causal one under the tiled protocol, and every
headline metric moves in the same direction (Table~\ref{tab:leak}). A trailing window tracks
amplitude drift within an epoch instead of imposing one scale across all 30 seconds, which
plausibly explains the gain. The important consequence is that neither arm depended on the leaked
information: removing it did not expose a hidden reliance, it removed a source of error. A pipeline
step that reads the future is therefore not buying accuracy in exchange for causality here, which
is worth stating plainly given how widely the step is used.

\begin{table}[t]
\caption{Effect of Closing the Preprocessing Leak (Sleep-EDF-78, C1, 15 Paired Runs Per Arm, Tiled
Evaluation)}
\label{tab:leak}
\centering
\footnotesize
\setlength{\tabcolsep}{4pt}
\begin{tabular}{lccc}
\toprule
 & Epoch $z$-score & Trailing window & $\Delta$ \\
\midrule
Causal $\kappa$         & 0.6406 & \textbf{0.6649} & $+0.0243$ \\
Non-causal $\kappa$     & 0.6693 & \textbf{0.6898} & $+0.0205$ \\
Causal accuracy         & 0.7289 & \textbf{0.7505} & $+0.0216$ \\
Causal macro-F1         & 0.6672 & \textbf{0.6894} & $+0.0222$ \\
Causal N1 F1            & 0.404  & \textbf{0.421}  & $+0.017$  \\
Causal $\kappa$ at 30~s & 0.6591 & \textbf{0.6847} & $+0.0256$ \\
\midrule
Pipeline reads the future & yes & \textbf{no} & --- \\
\bottomrule
\end{tabular}
\end{table}

\section{Model Architecture}

The network maps a window of raw, single-channel EEG to one five-class prediction per second in
four stages: a multi-resolution causal CNN (MRCNN), a causal TCN, causal masked self-attention, and
a linear classifier (Table~\ref{tab:tensors}).

\subsection{Multi-Resolution CNN}

Two causal convolution branches read the same raw signal in parallel: a 50-sample kernel (0.5~s at
100~Hz) targeting short events such as spindles and K-complex components, and a 400-sample kernel
(4~s) targeting slower activity such as slow waves. Each branch produces 16 channels; concatenation
gives 32 features per raw sample, and max-pooling with kernel and stride 100 reduces the 100~Hz
representation to one feature vector per second. Pooling is left-aligned, so only present and past
samples enter each pooled output.

\subsection{Temporal Convolutional Network}

Causality in each convolutional layer is enforced by left-padding only, so that an output depends
solely on present and past input:
\begin{equation}
y_t = \sum_{k=0}^{K-1} w_k \cdot x_{t-kd}
\end{equation}
where $K$ is the kernel size, $d$ the dilation, and the input is left-padded by $(K-1)\cdot d$
zeros so that $y_t$ never reads $x_i$ for $i > t$. Three residual blocks~\cite{resnet} of causal
dilated convolutions (dilations 1, 2, 4; kernel 3; two convolutions per block; 32 channels
throughout) follow this pattern, each with batch normalization~\cite{batchnorm}, ReLU, dropout
(0.2), and a residual connection. Dilation widens the causal receptive field without a larger
kernel: the three blocks give a 29-second causal receptive field on the per-second feature stream,
on top of the 4-second raw-signal window of the MRCNN's long branch.

Batch normalization pools statistics over the time axis during training, which is not a causal
operation. At inference the layers use frozen running statistics and are pure affine maps, so the
deployed model is causal and the perturbation test of Table~\ref{tab:verified} passes on it. We
note this because the training-time behaviour is easy to mistake for a violation of the claim.

\subsection{Causal Self-Attention and Classifier}

A four-head attention layer~\cite{attention}, followed by a residual connection and layer
normalization, lets each second attend to earlier seconds within the context window:
\begin{equation}
\mathrm{Attn}(Q,K,V) = \mathrm{softmax}\!\left(\frac{QK^{T}}{\sqrt{d_k}} + M\right)V
\end{equation}
where $M_{ij} = 0$ for $j \leq i$ and $M_{ij} = -\infty$ for $j > i$, so query position $i$ cannot
attend to key position $j$ in the future. Attention therefore extends the effective history to the
full context window (120~s in the primary configuration) while remaining strictly backward-looking.
A final linear layer maps the 32-dimensional representation at each second to five class logits;
softmax is not applied inside the model, since the training loss consumes logits directly.
Table~\ref{tab:params} gives the parameter budget.

\begin{table}[t]
\caption{Tensor Flow ($B$ = batch size, $L$ = context length in seconds)}
\label{tab:tensors}
\centering
\begin{tabular}{lll}
\toprule
Stage & Operation & Output shape \\
\midrule
Dataset    & Load $L$ one-second segments   & $(B, L, 100)$ \\
Input      & Flatten to one signal          & $(B, 1, L\cdot100)$ \\
MRCNN      & Short/long-scale causal convs. & $(B, 32, L)$ \\
TCN        & Causal dilated residual blocks & $(B, 32, L)$ \\
Attention  & Causal masked self-attention   & $(B, L, 32)$ \\
Classifier & Linear, 5 classes              & $(B, L, 5)$ \\
\bottomrule
\end{tabular}
\end{table}

\begin{table}[t]
\caption{Parameter Budget}
\label{tab:params}
\centering
\begin{tabular}{lr}
\toprule
Component & Parameters \\
\midrule
MRCNN branch, kernel 50 (0.5~s)   & 848 \\
MRCNN branch, kernel 400 (4~s)    & 6{,}448 \\
TCN residual block 1 (dilation 1) & 6{,}336 \\
TCN residual block 2 (dilation 2) & 6{,}336 \\
TCN residual block 3 (dilation 4) & 6{,}336 \\
Causal masked self-attention (4 heads) & 4{,}224 \\
Layer normalization               & 64 \\
Linear classifier ($32 \to 5$)    & 165 \\
\midrule
\textbf{Total} & \textbf{30{,}757} \\
\bottomrule
\end{tabular}
\end{table}

\subsection{Enforcing and Ablating Causality}

Causality is enforced in exactly two places: every causal convolution pads only the left side by
$(K-1)\cdot d$, per (3), and the attention mask $M$ in (4) zeros out any position after the query.
Setting the model to non-causal mode removes both, so that convolutions pad symmetrically and no
attention mask is applied, while every layer, channel width, and parameter remains identical. This
makes the comparison in Sections~\ref{sec:results-protocol} and \ref{sec:latency} a controlled one:
the two arms differ only in whether the network can access future signal, not in capacity.

\begin{table}[t]
\caption{Verified Properties (Repository Unit Tests)}
\label{tab:verified}
\centering
\footnotesize
\begin{tabular}{@{}p{0.56\columnwidth}l@{}}
\toprule
Property & Result \\
\midrule
Causal conv.\ preserves sequence length & PASS \\
Normalization ignores future samples & PASS (bit-identical) \\
Offline arrays equal the online filter & PASS ($5.2\times10^{-14}$) \\
Raw sample $\to$ normalization $\to$ logits & PASS (bit-identical) \\
Probe does perturb later outputs & PASS (non-vacuous) \\
Non-causal control does leak & PASS \\
Incremental buffer equals offline eval. & PASS (79{,}440~s, 0 diff.) \\
Normalizer is scale-invariant & PASS ($5\times10^{-7}$, $10^3\times$ gain) \\
Output shape and parameter count & PASS (30{,}757) \\
Hold-out splits disjoint and complete & PASS \\
Every subject tested once over 5 folds & PASS \\
Same seed reproduces the same split & PASS \\
\bottomrule
\end{tabular}
\end{table}

Table~\ref{tab:verified} reports the verification suite. Unlike the corresponding table in earlier
versions of this work, its scope is the whole path and not only the network: the perturbation test
runs from raw sample through normalization to logits, and a separate test streams a complete raw
recording one sample at a time through the deployment front end and requires the resulting labels
to match the offline evaluation exactly. Over one full held-out recording of 79{,}440 seconds there
were zero mismatches. The scale-invariance test was added for the second dataset, since DOD-H and
Sleep-EDF are recorded in different amplitude units, and confirms that the normalizer removes a
$1000\times$ change in input gain to within $5\times10^{-7}$.

\section{Experimental Setup}

\subsection{Configuration Registry}

Three architectural configurations recur throughout the ablations. \textbf{C0} is 60~s of context
with three TCN blocks. \textbf{C1} is 120~s of context with three TCN blocks, selected by the
context-length ablation (Section~\ref{sec:config-ablation}) and used for every Sleep-EDF-78 and
DOD-H result. \textbf{C2} is 120~s of context with four TCN blocks, the deeper variant ruled out by
the same ablation.

\subsection{Training}

Training uses subject-wise 5-fold cross-validation so that no subject's recordings cross the
train/validation/test boundary. The objective is focal loss~\cite{focal}
\begin{equation}
FL(p_t) = -\alpha_t (1-p_t)^{\gamma}\log(p_t)
\end{equation}
with focusing parameter $\gamma = 2.0$ and class weights $\alpha_t$ computed from the training
fold's own class counts, chosen over a weighted sampler for class imbalance, since combining both
pushed the model toward over-predicting N1 and reduced its precision to 0.12--0.20
(Section~\ref{sec:preproc-ablation}). Optimization uses a learning rate of 0.001, weight decay
$1\times10^{-4}$, batch size 64, early stopping (patience 8 epochs) on validation $\kappa$, and a
40-epoch cap; validation typically plateaus by epoch 10. Checkpoints are selected on validation
$\kappa$ and not macro-F1, since the two peak at different epochs. The held-out test fold is used
for neither parameter updates, early stopping, nor checkpoint selection.

For the causality comparison the full 5-fold protocol is repeated under three random seeds (42, 43,
44) for each arm on each dataset, giving 15 seed--fold result cells per arm per dataset, or 60
trained models in total.

\subsection{Evaluation Protocols}
\label{sec:protocol}

Two evaluation protocols appear in this paper and the distinction between them is the subject of
Section~\ref{sec:results-protocol}.

\textbf{Tiled.} The recording is covered by non-overlapping context windows and every position of
every window contributes one scored prediction. This is the default in every implementation we have
examined, including our own earlier work. Each window begins from a zero-padded empty history.

\textbf{Streaming.} The model recomputes every $s$ seconds against a full buffer and only the
freshest $s$ predictions of each window are kept, so after the opening window every scored second is
decided with at least $L_{\text{buf}} - s$ seconds of genuine history. This is what a deployed
system does. We use $s = 30$ throughout. Every second of every held-out recording is scored exactly
once under both protocols, so the two are directly comparable, and the same trained weights are used
for both: no model is retrained between protocols.

\subsection{Statistical Analysis}
\label{sec:stats}

The causality comparison is repeated under three random seeds. In the archived runs each seed
reseeds both the weight initialization and the subject partition, so the experiment is three
independent repeats of five-fold cross-validation over fifteen distinct train/test partitions drawn
from the subject pool, rather than five partitions measured three times. This is worth stating
precisely because it is easy to get wrong: fold 0 of seed 42 and fold 0 of seed 43 share 4 of their
16 Sleep-EDF test subjects. They are not the same fold.

Because repeats reuse subjects, the fifteen paired differences are not independent, and a naive
paired $t$-test over them overstates significance. The primary test is therefore the Nadeau--Bengio
corrected resampled $t$-test, which inflates the variance estimate by the train-set overlap between
repeats: for $k$-fold cross-validation over $n$ measurements the correction is $1/n + 1/(k-1)$,
here $1/15 + 1/4$. We report the naive statistic alongside it in the released analysis.

The pairing itself is exact. Within a seed, the causal and non-causal arms use byte-identical
subject lists, verified against the archived fold reports, so every one of the fifteen differences
is a properly matched pair. The protocol comparison of Section~\ref{sec:results-protocol} is paired
twice over: the same weights are scored under both protocols, so each arm's protocol gain is a
within-model difference, and the excess is the difference of two such gains on the same fold.

Confidence intervals on single-arm metrics are obtained by bootstrapping whole subjects rather than
seconds, with 2{,}000 draws. Seconds within one recording are heavily autocorrelated, so a
second-level interval would be far too narrow to mean anything.

\section{Results}

\subsection{Headline Results}

Table~\ref{tab:headline} reports the primary metrics for both datasets under both evaluation
protocols. The streaming protocol is the primary result, because it is what a deployed system
experiences. On Sleep-EDF-78 the model reaches $\kappa = 0.700$ and 77.7\% accuracy per second,
rising to $\kappa = 0.707$ and 78.2\% when predictions are majority-voted to 30~s epochs; on DOD-H
it reaches $\kappa = 0.667$ and 75.7\%.

\begin{table}[t]
\caption{Headline Results, Causal Arm, Subject-Wise 5-Fold Cross-Validation, Configuration C1, 15
Seed--Fold Cells Pooled Per Column}
\label{tab:headline}
\centering
\footnotesize
\setlength{\tabcolsep}{4pt}
\begin{tabular}{lcc|cc}
\toprule
 & \multicolumn{2}{c|}{Sleep-EDF-78} & \multicolumn{2}{c}{DOD-H} \\
Metric & Tiled & \textbf{Stream.} & Tiled & \textbf{Stream.} \\
\midrule
Accuracy (per second)         & 0.7505 & \textbf{0.7772} & 0.7319 & \textbf{0.7574} \\
Cohen's $\kappa$ (per second) & 0.6649 & \textbf{0.6998} & 0.6329 & \textbf{0.6674} \\
Macro-F1 (per second)         & 0.6894 & \textbf{0.7230} & 0.6765 & \textbf{0.7038} \\
N1 F1                         & 0.4212 & \textbf{0.4629} & 0.3515 & \textbf{0.3938} \\
Accuracy at 30~s              & 0.7657 & \textbf{0.7823} & 0.7526 & \textbf{0.7659} \\
$\kappa$ at 30~s              & 0.6847 & \textbf{0.7065} & 0.6603 & \textbf{0.6788} \\
Fold-to-fold $\kappa$ s.d.    & 0.0257 & 0.0277 & 0.0982 & 0.1063 \\
\midrule
Subjects / hours & \multicolumn{2}{c|}{78 / 1{,}629} & \multicolumn{2}{c}{25 / 205.5} \\
End-to-end causal pipeline & \multicolumn{2}{c|}{yes} & \multicolumn{2}{c}{yes} \\
Parameters & \multicolumn{4}{c}{30{,}757} \\
CPU inference & \multicolumn{4}{c}{0.026 ms/s} \\
\bottomrule
\end{tabular}
\end{table}

$\kappa$ is lower on Sleep-EDF-78 than on Sleep-EDF-20, which is expected: Sleep-EDF-78 spans ages
25--101 and is a harder population, and published baselines show the same direction, with AttnSleep
reporting 84.4\% on Sleep-EDF-20 against 81.6\% on Sleep-EDF-78. The fold-to-fold standard
deviation on DOD-H is four times that on Sleep-EDF-78, at 0.098 against 0.026, which is a direct
consequence of five held-out subjects per fold against sixteen and which governs what can and
cannot be resolved on that dataset (Section~\ref{sec:smallsample}).

\subsection{The Evaluation Protocol Accounts for Most of the Causality Penalty}
\label{sec:results-protocol}

Table~\ref{tab:protocol} is the central result of this paper. The same trained weights, scored under
the two protocols of Section~\ref{sec:protocol}, give different answers to the question of what
causality costs, and they do so on both datasets.

\begin{table}[t]
\caption{Cost of Causality Under Each Evaluation Protocol. 15 Paired Folds Per Row;
Nadeau--Bengio Corrected}
\label{tab:protocol}
\centering
\footnotesize
\setlength{\tabcolsep}{4pt}
\begin{tabular}{llccccc}
\toprule
Dataset & Protocol & Causal & Non-c. & $\Delta$ & Loses & $p$ \\
\midrule
Sleep-EDF-78 & Tiled     & 0.6649 & 0.6898 & $-0.0249$ & 14/15 & 0.005 \\
Sleep-EDF-78 & Streaming & \textbf{0.6998} & 0.6934 & $\mathbf{+0.0065}$ & 7/15 & 0.47 \\
\midrule
DOD-H        & Tiled     & 0.6329 & 0.6659 & $-0.0329$ & 10/15 & 0.22 \\
DOD-H        & Streaming & \textbf{0.6674} & 0.6711 & $\mathbf{-0.0037}$ & 6/15 & 0.89 \\
\bottomrule
\end{tabular}
\end{table}

Read on its own, Table~\ref{tab:protocol} is suggestive rather than decisive, and it is worth being
clear about why. On Sleep-EDF-78 the tiled penalty is resolved and the streaming one is not, which
is the shape the argument predicts. On DOD-H neither is resolved: with five subjects per held-out
fold the confidence interval on the tiled difference is $[-0.056, -0.012]$ and the corrected test
returns $p = 0.22$. A cost that cannot be resolved on one dataset cannot be shown to disappear on
it either, so a claim built on this table alone would rest on Sleep-EDF and be merely consistent
with DOD-H.

The decisive statistic is a different one, and it is decisive precisely because it does not require
the two arms to be comparable. Ask instead how much each arm gains from the protocol change, which
is a within-model difference on identical weights and identical held-out seconds
(Table~\ref{tab:asym}).

\begin{table}[t]
\caption{Protocol Gain by Arm: $\kappa$(Streaming) $-$ $\kappa$(Tiled), Paired Within Fold.
Nadeau--Bengio Corrected, $t(14)$}
\label{tab:asym}
\centering
\footnotesize
\setlength{\tabcolsep}{4pt}
\begin{tabular}{lccc}
\toprule
 & Sleep-EDF-78 & DOD-H & Both \\
\midrule
Causal arm gains       & $+0.0350$ & $+0.0345$ & --- \\
Non-causal arm gains   & $+0.0036$ & $+0.0052$ & --- \\
\midrule
\textbf{Excess to the causal arm} & $\mathbf{+0.0314}$ & $\mathbf{+0.0292}$ & --- \\
Folds with positive excess & 15/15 & 15/15 & \textbf{30/30} \\
Bootstrap 95\% CI & $[.0291,.0334]$ & $[.0253,.0332]$ & --- \\
Corrected $t(14)$ & $12.69$ & $6.48$ & --- \\
$p$ & $4.6\times10^{-9}$ & $1.4\times10^{-5}$ & --- \\
Per-seed range & $.0310$--$.0319$ & $.0285$--$.0297$ & --- \\
\bottomrule
\end{tabular}
\end{table}

The protocol change is worth roughly ten times as much to the causal arm as to the non-causal arm,
the excess is positive in every one of the thirty paired folds across the two datasets, and the
per-seed spread is under $0.001\,\kappa$ on each dataset. Two datasets that differ in electrode
montage, sampling rate, scoring team, class balance, subject count and recording length agree on
the size of the effect to within $0.002\,\kappa$. The lowest single-fold excess anywhere in the
thirty is $+0.0173$.

The aggregate causality cost follows from this arithmetic rather than being an independent
observation: the tiled penalty is $-0.0249$ and $-0.0329$, the excess is $+0.0314$ and $+0.0292$,
and the streaming penalty is what remains, $+0.0065$ and $-0.0037$. The protocol therefore accounts
for essentially all of the measured penalty on Sleep-EDF-78 and about 89\% of it on DOD-H.

The mechanism is direct. The streaming protocol keeps the freshest predictions of each window,
which carry 90--119 seconds of genuine past. Under tiling those same predictions are scored at every
buffer position, averaging roughly 60 seconds of history and reaching zero at the start of each
window. The causal model has nothing to work with there. The non-causal model at the same position
still has up to 119 seconds of lookahead, so it is barely affected.
Section~\ref{sec:latency} verifies this quantitatively rather than leaving it as an explanation.

\subsection{Accuracy Against Permitted Latency}
\label{sec:latency}

If the tiled score is an average over buffer positions, then measuring performance at each position
separately should both reproduce the tiled number and show where it comes from. It does, on both
datasets.

We define permitted latency $L$: the label for second $t$ is emitted at $t + L$, so the model sees
$L$ seconds of future and $L_{\text{buf}} - L$ seconds of past. $L = 0$ is real-time operation. A
single forward pass over a window yields a prediction at every position, and position $p$ carries
exactly $L_{\text{buf}} - 1 - p$ seconds of lookahead, so one sweep measures every latency. The
sweeps below use the seed-42 weights of each dataset with a 5-second window stride, scoring 1{,}171{,}020
Sleep-EDF-78 and 147{,}397 DOD-H held-out seconds at every latency.

Averaging the resulting curve over all 120 positions reproduces that model's tiled score to within
$0.003\,\kappa$ in all four cases: $0.6626$ against a measured $0.6634$ and $0.6890$ against
$0.6912$ for the two Sleep-EDF arms, and $0.6537$ against $0.6567$ and $0.6915$ against $0.6922$ for
the two DOD-H arms. The tiled protocol is, arithmetically, this average.

Table~\ref{tab:quartile} shows where the penalty lives, and the two datasets show the same shape.
At buffer positions with 91--120 seconds of history the two arms are level. The entire aggregate
penalty is generated by positions with under roughly 60 seconds of history, and the majority of it
by positions with under 30: across the last quartile the causal arm loses $0.114$ and $0.128\,\kappa$
against its own best quartile, while the non-causal arm loses $0.018$ and $0.028$.

\begin{table}[t]
\caption{Where the Tiled Penalty Comes From. $\kappa$ by Buffer Quartile, Seed 42, 120~s Buffer.
Positive $\Delta$ Means the Causal Arm Pays}
\label{tab:quartile}
\centering
\footnotesize
\setlength{\tabcolsep}{3.6pt}
\begin{tabular}{lccc|ccc}
\toprule
 & \multicolumn{3}{c|}{Sleep-EDF-78} & \multicolumn{3}{c}{DOD-H} \\
History (s) & Caus. & Non-c. & $\Delta$ & Caus. & Non-c. & $\Delta$ \\
\midrule
91--120 & \textbf{0.6982} & 0.6930 & $-0.0052$ & 0.6904 & 0.6961 & $+0.0057$ \\
61--90  & 0.6923 & 0.6969 & $+0.0046$ & 0.6887 & 0.7044 & $+0.0157$ \\
31--60  & 0.6753 & 0.6911 & $+0.0158$ & 0.6737 & 0.6970 & $+0.0233$ \\
1--30   & 0.5847 & 0.6752 & $+0.0905$ & 0.5622 & 0.6685 & $+0.1063$ \\
\bottomrule
\end{tabular}
\end{table}

Fig.~\ref{fig:latency} gives the full curves. The causal arm declines monotonically from $L = 0$ on
both datasets: it cannot read future signal, so a larger $L$ only costs it history, and at
$L = 119$, with one second of history, it collapses to $\kappa = 0.317$ and $0.252$. The non-causal
arm traces an inverted U, gaining lookahead and losing history, and peaks at $L = 38$~s on Sleep-EDF-78
and $L = 46$~s on DOD-H.

The comparison that matters is between each arm's best achievable operating point, and here the two
datasets disagree. On Sleep-EDF-78 the causal model achieves $\kappa = 0.7008$ at zero latency and
the non-causal model's optimum is $\kappa = 0.6977$, at 38 seconds of delay: lookahead buys nothing
that history does not already provide, though the margin is $0.003\,\kappa$ and we would not defend
it as more than a tie. On DOD-H the causal model achieves $\kappa = 0.6922$ at zero latency and the
non-causal optimum is $\kappa = 0.7058$, at 46 seconds, exceeding it by $0.0136$. The crossover,
the latency beyond which the non-causal arm leads, is 24 seconds on Sleep-EDF-78 and 5 seconds on
DOD-H.

We report this disagreement rather than the agreement because it bounds the claim. What replicates
across the two datasets is the protocol asymmetry, at $+0.0314$ and $+0.0292\,\kappa$ with 30 of 30
folds. What does not replicate is the stronger statement that a non-causal model cannot beat a
causal one at any latency: that holds on Sleep-EDF-78 and fails on DOD-H, where roughly
$0.014\,\kappa$ is available in exchange for 46 seconds of delay. A designer choosing between the
two should read the first result as general and the second as dataset-specific.

\begin{figure}[t]
\centering
\includegraphics[width=\columnwidth]{figures/fig_latency.png}
\caption{Cohen's $\kappa$ against permitted latency $L$ on both datasets, both arms, 5-point moving
average over $L$. Larger $L$ means less history and more lookahead within a fixed 120~s buffer. The
causal arm is best at $L = 0$ and declines monotonically; the non-causal arm peaks at $L = 38$ and
$46$~s. The causal collapse at large $L$ is the cold-start behaviour that the tiled protocol
averages into its score.}
\label{fig:latency}
\end{figure}

There is a reading of the $L = 0$ comparison that we find persuasive but cannot prove from these
experiments alone. At zero lookahead the non-causal architecture has no advantage left to use, but
it still carries symmetric padding and an unmasked attention, both of which spend capacity on
positions that contain nothing. Under that reading the causal constraint stops being a handicap and
becomes a useful prior when future signal is unavailable at inference. Confirming it would require
varying the constraint independently of the deployment protocol, which we have not done.

\subsection{Context Length and the Fixed-Buffer Objection}
\label{sec:buffer}

Section~\ref{sec:latency} measures latency within a buffer of fixed length, so permitting $L$
seconds of lookahead necessarily costs $L$ seconds of history. An offline system would not accept
that trade: it would enlarge the buffer and take both. Whether the result survives a larger buffer
is therefore a fair question, and we address it as far as these weights allow. This subsection uses
Sleep-EDF-78 only.

For the causal arm the question is answered exactly. Left-padded convolutions, left-aligned pooling
and a causal attention mask together make the output at buffer position $p$ a function of inputs
$0 \ldots p$ only, so scoring position $p$ of a long window is bit-identical to scoring the last
position of a window of length $p+1$. Evaluating one long buffer therefore yields the model's entire
context-response curve at once, and that curve cannot depend on the buffer it was measured in.
Table~\ref{tab:context} confirms this empirically: measured in buffers of 120, 240 and 360~s,
$\kappa$ at any fixed history agrees to within $0.0035$.

\begin{table}[t]
\caption{Causal Arm: $\kappa$ Against Available History (Sleep-EDF-78, C1, 360~s Buffer)}
\label{tab:context}
\centering
\footnotesize
\setlength{\tabcolsep}{3.4pt}
\begin{tabular}{lccccccc}
\toprule
History (s) & 30 & 60 & 90 & 120 & \textbf{145} & 240 & 360 \\
\midrule
$\kappa$ & 0.658 & 0.690 & 0.701 & 0.704 & \textbf{0.705} & 0.698 & 0.683 \\
\bottomrule
\end{tabular}
\end{table}

The curve rises steeply to about 90~s, is nearly flat from 90 to 180~s, peaks at approximately
145~s, and declines beyond. The optimum is located independently at 145~s in the 240~s buffer and
146~s in the 360~s buffer. Two things follow. A deployed causal system should use a buffer of
roughly 145~s rather than the 120~s used here, which is worth about $0.005\,\kappa$ at no latency
cost. And the decline past 180~s is at least partly a train/test length mismatch, since these
weights were trained on 120~s sequences; it does not establish that a model trained at longer
context would not benefit.

For the non-causal arm the same experiment is not interpretable, and it is worth being explicit
about why. Enlarging the buffer necessarily lengthens the sequence its unmasked attention operates
over, and unlike the causal arm it has no mask limiting the effective span.
Table~\ref{tab:contamination} isolates the effect: holding history fixed and increasing only the
available lookahead reduces $\kappa$ monotonically.

\begin{table}[t]
\caption{Non-Causal Arm: $\kappa$ at Fixed History as the Buffer Grows (Sleep-EDF-78)}
\label{tab:contamination}
\centering
\begin{tabular}{lccc}
\toprule
 & \multicolumn{3}{c}{Buffer} \\
History & 120~s & 240~s & 360~s \\
\midrule
\hphantom{0}60~s & 0.6942 & 0.6714 & 0.6550 \\
\hphantom{0}90~s & 0.6971 & 0.6820 & 0.6666 \\
120~s            & 0.6870 & 0.6886 & 0.6757 \\
\bottomrule
\end{tabular}
\end{table}

Additional lookahead cannot itself degrade a model that is free to ignore it, so this is the
signature of attention operating over sequences longer than those it was trained on rather than a
property of lookahead. Consequently the natural test, giving the non-causal arm 120~s of history and
120~s of lookahead in a 240~s buffer, returns $\kappa = 0.689$, which is below the causal arm's
$0.704$ at the same history with no lookahead, but which we do not treat as decisive: the figure is
depressed by an amount these experiments cannot separate from the comparison.

We therefore state the scope of Section~\ref{sec:latency}'s result precisely. Within the buffer both
arms were trained for, the causal model at zero latency is not beaten by the non-causal model at any
latency on Sleep-EDF-78, and is beaten by $0.014\,\kappa$ at 46~s on DOD-H. Whether a non-causal
model \emph{trained} at a longer context, and so able to use full history and full lookahead
together, would exceed it is untested here and is the clearest experiment this analysis calls for.
It requires retraining both arms at the longer context and is left to future work. Note that the
protocol asymmetry of Section~\ref{sec:results-protocol} is unaffected by this objection: it is a
within-arm quantity measured at a single buffer length, so it makes no comparison between arms that
a longer buffer could overturn.

\subsection{Per-Class Cost}
\label{sec:perclass}

Table~\ref{tab:perclass} applies the same comparison per class. Under the tiled protocol REM appears
to pay $0.060$ F1 on both datasets, and on Sleep-EDF-78 it is the only class whose penalty survives
Holm correction across the five per-class tests. That result does not survive a change of protocol.
Under the streaming protocol the REM difference is $+0.002$ on Sleep-EDF-78 and $+0.023$ on DOD-H,
and N1, N2 and N3 invert on both datasets, the causal arm ending ahead on all three.

\begin{table}[t]
\caption{Per-Class F1 by Evaluation Protocol, 15 Paired Folds Per Column. Positive $\Delta$ Means
the Causal Arm Pays}
\label{tab:perclass}
\centering
\footnotesize
\setlength{\tabcolsep}{3.2pt}
\begin{tabular}{lcc|cc}
\toprule
 & \multicolumn{2}{c|}{Sleep-EDF-78} & \multicolumn{2}{c}{DOD-H} \\
Class & Tiled $\Delta$ & Stream.\ $\Delta$ & Tiled $\Delta$ & Stream.\ $\Delta$ \\
\midrule
W   & $+0.018$ & $+0.006$ & $+0.046$ & $+0.035$ \\
N1  & $+0.025$ & $-0.009$ & $+0.007$ & $-0.018$ \\
N2  & $+0.002$ & $-0.017$ & $+0.014$ & $-0.004$ \\
N3  & $+0.008$ & $-0.020$ & $+0.015$ & $-0.004$ \\
REM & $\mathbf{+0.060}$ & $\mathbf{+0.002}$ & $\mathbf{+0.060}$ & $\mathbf{+0.023}$ \\
\midrule
Macro & $+0.023$ & $-0.007$ & $+0.028$ & $+0.006$ \\
\bottomrule
\end{tabular}
\end{table}

The explanation is the same as for the aggregate, and it is mechanistically sensible. REM is defined
partly by the context surrounding it, which is why a technician resolves REM against neighbouring
epochs rather than scoring it in isolation. It is therefore the class most damaged by a nearly empty
buffer, and the cold-start positions that the tiled protocol averages in are exactly where that
damage occurs. Fig.~\ref{fig:perclass} shows the two readings side by side.

Two residuals are worth naming rather than smoothing over. REM still costs the causal arm $0.023$ F1
on DOD-H under the deployment protocol, roughly ten times what it costs on Sleep-EDF-78, and Wake
costs $0.035$ on DOD-H under both protocols and almost nothing on Sleep-EDF-78. The Wake difference
is plausibly a consequence of the dataset rather than of causality: DOD-H is not wake-trimmed, so its
Wake is sleep-adjacent and transitional throughout, and transitional Wake is exactly the case where
seeing forward helps. We offer that as a reading, not a measurement.

What survives is a weaker and more careful statement than the earlier one: under tiled evaluation
the apparent causality penalty is concentrated in REM, and this is because REM is the class most
sensitive to history, not because REM depends on the future in a way the causal constraint forbids.

\begin{figure}[t]
\centering
\includegraphics[width=\columnwidth]{figures/fig_perclass.png}
\caption{Per-class F1 the causal arm gives up, under batch tiling and at deployment latency, on both
datasets. The apparent REM penalty shrinks by an order of magnitude on Sleep-EDF-78 and by a factor
of about three on DOD-H, while N1, N2 and N3 invert on both.}
\label{fig:perclass}
\end{figure}

\subsection{Temporal Stability and Causal Smoothing}

Because the model labels every second independently, its raw output is markedly more fragmented than
a scored hypnogram. On the Sleep-EDF-78 held-out recordings it changes stage 167 times an hour
against 13 times an hour for the technician. A trailing-window mode filter, in which the label at
second $t$ becomes the most common prediction over the preceding $w$ seconds,
\begin{equation}
\hat{y}_t = \mathrm{mode}(\{\hat{y}'_{t-w+1},\dots,\hat{y}'_t\})
\end{equation}
uses only past predictions, preserves causality, and requires no retraining
(Table~\ref{tab:smoothing}).

\begin{table}[t]
\caption{Causal Mode-Filter Smoothing (Sleep-EDF-78, C1, Streaming Pipeline, Seed 42)}
\label{tab:smoothing}
\centering
\begin{tabular}{lccccr}
\toprule
Window (s) & Acc. & $\kappa$ & M-F1 & N1 F1 & Chg./hr \\
\midrule
1 (raw) & 0.7487 & 0.6634 & 0.6903 & 0.414 & 167 \\
5       & 0.7514 & 0.6668 & 0.6931 & 0.416 & 68 \\
10      & 0.7541 & 0.6702 & 0.6960 & \textbf{0.418} & 53 \\
15      & 0.7548 & 0.6711 & 0.6969 & 0.417 & 35 \\
\textbf{30} & \textbf{0.7568} & \textbf{0.6736} & \textbf{0.6993} & 0.417 & 25 \\
45      & 0.7567 & 0.6732 & 0.6993 & 0.414 & 18 \\
60      & 0.7568 & 0.6732 & 0.6995 & 0.411 & 15 \\
90      & 0.7549 & 0.6703 & 0.6970 & 0.401 & 11 \\
120     & 0.7522 & 0.6663 & 0.6937 & 0.393 & 9 \\
\bottomrule
\end{tabular}
\end{table}

A 30~s window is optimal for $\kappa$ ($+0.0102$ over raw) and reduces fragmentation sevenfold; a
45~s window brings the stage-change rate to 18 per hour, close to the human rate, at negligible
cost. N1 F1 peaks earlier, around a 10~s window, and declines as smoothing increases, because N1
bouts are short and heavy smoothing absorbs them into neighbouring stages. The optimal window
therefore depends on whether overall agreement or transitional-stage sensitivity is prioritized. One
caveat applies: the smoothing window was selected on the same held-out folds it is reported on, so
the $+0.0102$ figure is a best-case estimate; selecting $w$ on validation data would be the stricter
protocol. This sweep was run on Sleep-EDF-78 only.

\subsection{What the Per-Second Output Rate Buys}
\label{sec:boundary}

A per-second output rate is only useful if it lets a system report something sooner. We compare when
each system can report a stage change at all, under the causal constraint both face in deployment:
the per-second system emits the label for second $t$ at second $t$, whereas an epoch-level system
cannot emit the label for epoch $k$ until epoch $k$ has finished, at second $30(k+1)$. The epoch
system therefore carries a structural 30-second floor on any change that begins an epoch.

Both systems are charged for their evidence: a detection requires the new stage to be held for $h$
seconds, so the per-second system reports at the last second of that run rather than the first.
Changes the model was already calling during the run-up are excluded, since those are false
positives during the outgoing stage rather than early detections, and label changes at a night join
are skipped.

At $h = 10$~s, on the conservative subset, the per-second system reports the change first in 94.8\%
of transitions on Sleep-EDF-78 and 95.6\% on DOD-H, with a median 19 seconds saved on both. On
Sleep-EDF-78 the advantage holds across all five stages entered (+19 to +34 seconds).

This result is sensitive to $h$ and it is important to say so. Sweeping $h$ from 3 to 30 seconds on
Sleep-EDF-78, the share of transitions the per-second system reports first falls from 100\% at
$h = 3$ to 94.8\% at $h = 10$, 66.5\% at $h = 20$ and 29.6\% at $h = 30$, where the epoch system
leads. This is mechanically expected: requiring 30 seconds of confirmation reimposes the epoch and
discards the resolution. The honest statement is therefore conditional. A finer output rate is worth
having precisely when the application is willing to act on less than an epoch of evidence, and is
worth nothing when it is not. We report the full sweep rather than a single figure, and note that
the scored transition boundary is itself only located to the nearest 30~s, so these numbers describe
latency against the scored transition and not against the physiological one.

\subsection{Configuration Ablation}
\label{sec:config-ablation}

On Sleep-EDF-20, doubling context from 60~s (C0) to 120~s (C1) improved $\kappa$ by $+0.0192$;
adding a fourth dilated block (C2) changed $\kappa$ by $-0.0050$, within fold variance, so C1 was
retained for the Sleep-EDF-78 and DOD-H runs (Table~\ref{tab:config}). Sequence length changes
computation and available history but not parameter count; the fourth block adds parameters without
benefit.

\begin{table}[t]
\caption{Configuration Ablation, Sleep-EDF-20, Tiled Evaluation}
\label{tab:config}
\centering
\begin{tabular}{lcccc}
\toprule
Config. & Acc. & $\kappa$ & M-F1 & N1 F1 \\
\midrule
C0 (60~s, 3 blocks)  & 0.7297 & 0.6433 & 0.6713 & 0.324 \\
C1 (120~s, 3 blocks) & 0.7436 & 0.6625 & 0.6879 & 0.336 \\
C2 (120~s, 4 blocks) & 0.7405 & 0.6575 & 0.6851 & 0.348 \\
\bottomrule
\end{tabular}
\end{table}

Given the latency result, the context-length finding reads differently than it did before. If
history substitutes for lookahead, then more history should continue to help a causal model, and
60~s to 120~s buying $0.019\,\kappa$ is consistent with that. Section~\ref{sec:buffer} extends the
curve post hoc and locates its optimum at about 145~s.

\subsection{The Small-Sample Warning, Twice}
\label{sec:smallsample}

Under the tiled protocol on Sleep-EDF-78 the causality effect is $-0.0249\,\kappa$ with the causal
arm losing in 14 of 15 paired runs, corrected $t(14) = -3.36$, $p = 0.005$. The same experiment on
the 20-subject split yields $+0.0143\,\kappa$ with $t(4) = 1.26$, $p = 0.276$: the wrong sign, and
not significant. Fold-to-fold variance at 20 subjects is 3.1 times larger than at 78 ($\kappa$
s.d.\ 0.092 against 0.030, both from the epoch-normalised runs on which that experiment was
performed), enough to obscure an effect of this magnitude and even invert it.

DOD-H repeats the lesson independently and at a scale typical of this literature. With 25 subjects
and five per held-out fold, the fold-to-fold $\kappa$ standard deviation is $0.098$, four times that
of Sleep-EDF-78, and the corrected test on the tiled causality difference returns $p = 0.22$ despite
a point estimate ($-0.0329$) larger than Sleep-EDF-78's. The same fifteen runs resolve the protocol
asymmetry at $p = 1.4\times10^{-5}$, because that quantity is a within-model difference whose
between-fold variance is an order of magnitude smaller. The general point is not that DOD-H is too
small to study, but that at this subject count the choice of estimand decides what is measurable:
between-arm differences are not resolvable and within-arm differences are.

We therefore retain and strengthen the methodological warning. Causality penalties measured on
20-to-25-subject splits, common in this literature for compute reasons, warrant caution, because the
true effect can be undetectable or sign-reversed from variance alone.

\subsection{Preprocessing Ablation}
\label{sec:preproc-ablation}

Table~\ref{tab:preproc} isolates the effect of wake trimming and the weighted sampler on
Sleep-EDF-20. Without trimming, Wake constitutes 68\% of the data and inflates accuracy; combining a
weighted sampler with focal loss additionally reduced N1 precision to 0.12--0.20. Correcting both
improved every individual sleep-stage F1, while overall accuracy decreased because the easy Wake
majority previously carrying it was removed. This is a clear case where headline accuracy moves
opposite to model quality.

\begin{table}[t]
\caption{Preprocessing Ablation, Sleep-EDF-20}
\label{tab:preproc}
\centering
\begin{tabular}{lccc}
\toprule
Metric & Untrim., sampler on & Trim., sampler off & $\Delta$ \\
\midrule
Accuracy & 0.8064 & 0.7283 & $-0.0781$ \\
Kappa    & 0.6584 & 0.6427 & $-0.0157$ \\
Macro-F1 & 0.6322 & 0.6786 & $+0.0463$ \\
N1 F1    & 0.2481 & 0.3554 & $+0.1073$ \\
N2 F1    & 0.7045 & 0.7829 & $+0.0784$ \\
N3 F1    & 0.6890 & 0.7752 & $+0.0863$ \\
REM F1   & 0.5898 & 0.6784 & $+0.0886$ \\
Wake F1  & 0.9296 & 0.8010 & $-0.1286$ \\
\bottomrule
\end{tabular}
\end{table}

\subsection{Comparison with Published Baselines}
\label{sec:baselines}

Table~\ref{tab:baselines} compares against published single-channel baselines on Sleep-EDF-20, the
split on which those baselines report. This model's rows are its Sleep-EDF-20 results (C1), so the
comparison is on a matched subject pool.

DeepSleepNet and AttnSleep remain ahead on both metrics. Both read future signal, and DeepSleepNet
uses roughly 700 times more parameters. What this model offers against them is not higher accuracy
but a 30$\times$ finer output rate at roughly 1/700th the parameter count, with the causal property
established rather than assumed, and now with the accompanying result that the causal constraint is
not what separates them. Section~\ref{sec:latency} bounds the contribution of causality to that gap
at approximately zero under matched deployment conditions, which leaves the remainder to capacity,
architecture and training scale.

\begin{table}[t]
\caption{Comparison with Published Baselines, Sleep-EDF-20}
\label{tab:baselines}
\centering
\begin{tabular}{lccccc}
\toprule
Model & Acc. & $\kappa$ & Params & Rate & Causal \\
\midrule
DeepSleepNet~\cite{deepsleepnet} & 82.2\% & 0.754 & $\sim$22M & 30~s & no \\
AttnSleep~\cite{attnsleep}$^{a}$ & 80.9\% & 0.709 & $\sim$48K & 30~s & partial \\
This work, 30~s agg.             & 76.1\% & 0.684 & 30.7K & 30~s & yes$^{b}$ \\
This work, per-second            & 74.4\% & 0.662 & 30.7K & \textbf{1~s} & yes$^{b}$ \\
\bottomrule
\end{tabular}
\begin{flushleft}
\footnotesize
$^{a}$ In-house rerun at AttnSleep's native 30-second output over the same 5 folds used throughout
this table (80.90\%, $\kappa$ 0.7088). AttnSleep's own report quotes 84.4\% on Sleep-EDF-20 and
81.6\% on Sleep-EDF-78 under its own protocol; those figures are not interchangeable with this
table's.\\
$^{b}$ Causal through the whole path, verified by the perturbation and streaming-equivalence tests of
Table~\ref{tab:verified}.
\end{flushleft}
\end{table}

CareSleepNet~\cite{caresleepnet} is discussed rather than tabulated. Its frequently quoted figure of
78.07\%/0.694 is its best fold, not its mean; its own report shows fold 1 at 38.12\% and fold 4 as
unrecoverable. Without access to its complete per-fold results a fair mean-based comparison is not
possible, so no head-to-head claim is made here. We likewise do not tabulate DOD-H against published
DOD-H results, since those use the full multi-channel montage and a different scoring consensus, and
the DOD-H experiment here exists to test a protocol effect rather than to compete on that
leaderboard.

\subsection{Efficiency}

Latency was benchmarked on an Intel Core i9-14900HX CPU using a corrected protocol: batch size 1,
single thread, 100 warm-up and 1{,}000 timed runs, median and interquartile range reported, with the
normalizer measured separately. C1 costs $0.0271$~ms per second of EEG and C0 $0.0216$~ms.

This metric is per-output throughput, not per-invocation latency: it is total wall-clock time for a
synthetic sequence divided by its length in seconds of EEG. A device labelling every second pays a
rolling cost instead, which is roughly forty times larger and is what
\texttt{scripts/streaming\_demo.py --mode rolling} measures. Absolute per-invocation latency, peak
memory, and MAC counts are not measured here and are needed before any claim about a specific target
device.

\section{Discussion}

The first thing to say is what changed and why. An earlier version of this work reported that the
causal constraint costs $0.029\,\kappa$ and that the cost concentrates in REM. Both statements are
reproducible under the evaluation protocol they were measured with, and both are largely artifacts
of that protocol. We report them here alongside the corrected measurement rather than quietly
replacing them, because the discrepancy is the finding.

The mechanism is not subtle once seen. A streaming model is defined by what it can do with the
history it has accumulated. Tiling a recording with independent context windows repeatedly destroys
that history and then scores the model on the wreckage, and it does so asymmetrically: a model that
can read forward is largely unaffected at the start of a window, and a model that cannot is crippled
there. The tiled score is the mean over buffer positions, which we verify to within $0.003\,\kappa$
on four independent model--dataset combinations, and the causality penalty is generated almost
entirely by positions a deployed system occupies once and then leaves.

The reason we lead with the asymmetry rather than with the disappearance of the penalty is
statistical, and it is the methodological point we would most like to be adopted. A between-arm
difference on 25 subjects is not resolvable; the corrected test returns $p = 0.22$ on DOD-H even
though the point estimate is larger than Sleep-EDF-78's. The same fifteen runs resolve the
within-arm protocol gain at $p = 1.4\times10^{-5}$, because scoring identical weights on identical
seconds under two protocols removes model variance from the comparison entirely. Any group with a
20-to-30-subject corpus can run this measurement and get a decisive answer, where the between-arm
ablation they would naturally reach for cannot give them one.

We do not claim from this that causality is free in general. Three things bound the claim. First,
the streaming protocol denies the non-causal arm the lookahead that defines it, so this is a
comparison of deployment options and not a clean architectural ablation; the fair framing is matched
latency, and at matched latency the answer differs by dataset, with lookahead worth nothing on
Sleep-EDF-78 and about $0.014\,\kappa$ at 46 seconds of delay on DOD-H. Given unbounded latency,
which is the offline setting, lookahead is worth roughly $0.025$ to $0.033\,\kappa$ under tiling.
Second, the non-causal arm here is our own matched twin rather than a state-of-the-art offline
model, so the comparison bounds what causality costs within this capacity class and not in general.
Third, the buffer is fixed at the length both arms were trained for; Section~\ref{sec:buffer}
explains what that does and does not permit, and notes that the asymmetry result is a within-arm
quantity and so untouched by the objection.

The practical implication is narrow and useful. A system that must operate in real time should not
be evaluated with a protocol that makes it look worse than it is, and reported causality penalties
in this literature should be read with the evaluation protocol in mind. We would encourage any
comparable ablation to report both protocols, since on our data the two numbers differ by more than
the effect being measured.

The preprocessing finding deserves emphasis for the same reason. A network can be causal in every
layer, verified bit-for-bit, and still sit downstream of a normalization step that reads the future.
That configuration is easy to build accidentally and hard to notice, because the network passes
every causality test one would think to write. We found it by verifying our own work rather than
because a reviewer asked, and the verification in Table~\ref{tab:verified} is built to cover the
whole path for exactly that reason.

Finally, the per-second output rate. Section~\ref{sec:boundary} gives the first quantitative answer
we are aware of to what a finer rate buys: at a 10-second confirmation window, transitions are
reported a median 19 seconds sooner than an epoch-level system can report them, in 94.8\% and 95.6\%
of cases on the two datasets. That advantage is conditional on the confirmation window and
disappears by 30 seconds, which is the honest form of the claim and also the useful one: it tells a
system designer exactly when the finer rate is worth the fragmentation it brings.

\section{Limitations}

This work covers two datasets, two electrode derivations and one architecture. Generalization to
other sensors, to clinical populations, and above all to a second architecture is untested, and the
protocol effect is a claim about an interaction between a model class and an evaluation procedure,
so its magnitude on a recurrent or encoder--decoder model is unknown. The mechanism argues that any
model whose predictions depend on accumulated history should show it, and neither dataset here
contradicts that, but two instances of one architecture is not a general result.

The two datasets are also not independent in every respect: both are single-channel EEG scored to
AASM conventions at 30~s resolution, and DOD-H was chosen partly because the pipeline could be
ported to it without changing the network. A corpus with different recording hardware and a
different scoring tradition would be a stronger test.

Supervision is 30-second expert labels replicated to 1~Hz, because no public dataset currently
provides per-second ground truth, so per-second accuracy should be read as an approximation and
$\kappa$ treated as the more meaningful summary. Section~\ref{sec:boundary} measures when a
per-second system can report a change rather than how accurately it localizes one, precisely because
the scored boundary is itself only known to 30 seconds.

No human inter-scorer agreement ceiling is used here. DOD-H is scored by five technicians and could
in principle supply one; we use the released consensus hypnogram only, and quantifying the model
against inter-scorer spread on that dataset is left undone.

Two free parameters are untested. The normalization window was fixed at 30~s to match the epoch
length it replaces, and no other setting was trained. The context length was fixed at 120~s;
Section~\ref{sec:buffer} shows post hoc that approximately 145~s would have been better for the
causal arm, and that the question of whether a non-causal model trained at longer context could use
history and lookahead together to exceed the causal arm at zero latency is open. Both require
retraining and neither is resolved here.

The causality estimate rests on fifteen paired runs per dataset, and those runs reuse subjects across
repeats, which is what the Nadeau--Bengio correction accounts for. A design with the partition held
fixed across seeds would separate initialization variance from partition variance and we did not have
one. On DOD-H specifically, five subjects per held-out fold is too few to resolve the between-arm
difference at all, which we state as a limitation of that dataset rather than as a result.

The smoothing sweep and the fixed-buffer analysis were run on Sleep-EDF-78 only. The smoothing window
was selected on the same held-out folds used to report it, so its gain is a best-case estimate.
Efficiency is characterized by CPU throughput alone; MAC counts, peak inference memory,
per-invocation latency, and measurements on an actual wearable-class device are not reported, so the
real-time claim should be read as evidence of feasibility on a desktop CPU, not as a device-level
result.

\section{Conclusion}

A strictly causal, per-second sleep-staging model was built, verified end-to-end from raw sample to
logit, and compared against a non-causal twin identical in every other respect, on 78 Sleep-EDF
subjects and on the 25 subjects of DOD-H. Measured with the batch-tiling protocol standard in this
literature, the causal constraint costs $0.025$ and $0.033\,\kappa$ and appears on both datasets to
fall on REM. Most of that is produced by the protocol. Moving the same weights to a streaming
protocol lifts the causal arm roughly ten times as much as the non-causal arm, by an excess of
$+0.0314$ and $+0.0292\,\kappa$ that is positive in all 30 paired folds and resolved at
$p = 4.6\times10^{-9}$ and $p = 1.4\times10^{-5}$. The tiled score is arithmetically the mean over
context-buffer positions, verified to within $0.003\,\kappa$ in four model--dataset combinations, and
the penalty it reports is generated by positions a streaming system occupies once at start-up. Once
that artifact is removed the residual cost is $+0.0065$ and $-0.0037\,\kappa$, neither resolvable,
and the apparent REM concentration falls from $0.060$ F1 to $0.002$ and $0.023$.

Two further results support the same theme, that causality in this setting is a property of the
pipeline and the protocol rather than only of the network. A look-ahead leak in per-epoch
normalization, near-universal in this literature, was identified and closed, and closing it improved
both arms. And a parameter-free causal smoothing step recovers output stability at inference with no
retraining.

At 30{,}757 parameters and $0.026$~ms per second of EEG the model is compact and fast enough for
on-device operation while emitting labels 30$\times$ more often than epoch-level systems. Whether the
protocol effect reproduces on a second architecture, and whether a non-causal model trained at longer
context can use history and lookahead together to beat a causal one at zero latency, are the two
clearest directions this measurement opens.

\section*{Reproducibility and Data Availability}
\sloppy

Subject-wise 5-fold splits, per-fold checkpoints, archived per-fold result files, the statistical
analysis, the latency sweeps, the smoothing sweep, and the boundary-latency analysis for both
datasets are available at \url{https://github.com/swapnil5053/causal-sleep-staging}, alongside the
implementation, configuration files, evaluation scripts, the DOD-H preprocessing script and its
manifest, and unit tests covering the causality and split properties reported in
Table~\ref{tab:verified}. Per-second predictions and logits are archived for every held-out fold, so
every analysis in Sections~\ref{sec:results-protocol}--\ref{sec:boundary} can be reproduced without
retraining. A camera-ready submission will cite an archived snapshot with a persistent identifier and
a pinned commit hash rather than a branch reference.

The Sleep-EDF Expanded dataset is publicly available from PhysioNet~\cite{physionet,kemp} under the
Open Data Commons Attribution License, and DOD-H is publicly available from the Dreem Open
Datasets~\cite{dod}; neither is redistributed here. This work is a secondary analysis of
de-identified public data and involved no new human-subject data collection, so institutional ethics
review was not required.

\section*{Acknowledgment}

The authors thank Dr.\ Bhaskarjyoti Das for supervision and guidance throughout this project.

\section*{Declaration of Competing Interests}

The authors declare no competing financial or non-financial interests. This research received no
specific grant from any funding agency in the public, commercial, or not-for-profit sectors.

\begin{thebibliography}{99}

\bibitem{deepsleepnet} A. Supratak, H. Dong, C. Wu, and Y. Guo, ``DeepSleepNet: A model for automatic
sleep stage scoring based on raw single-channel EEG,'' \emph{IEEE Trans. Neural Syst. Rehabil. Eng.},
vol. 25, no. 11, pp. 1998--2008, Nov. 2017.

\bibitem{attnsleep} E. Eldele, Z. Chen, C. Liu, M. Wu, C.-K. Kwoh, X. Li, and C. Guan, ``An
attention-based deep learning approach for sleep stage classification with single-channel EEG,''
\emph{IEEE Trans. Neural Syst. Rehabil. Eng.}, vol. 29, pp. 809--818, 2021.

\bibitem{caresleepnet} J. Wang, S. Zhao, H. Jiang, Y. Zhou, Z. Yu, T. Li, S. Li, and G. Pan,
``CareSleepNet: A hybrid deep learning network for automatic sleep staging,'' \emph{IEEE J. Biomed.
Health Inform.}, vol. 28, no. 12, pp. 7392--7405, Dec. 2024.

\bibitem{kemp} B. Kemp, A. H. Zwinderman, B. Tuk, H. A. C. Kamphuisen, and J. J. L. Oberye,
``Analysis of a sleep-dependent neuronal feedback loop: the slow-wave microcontinuity of the EEG,''
\emph{IEEE Trans. Biomed. Eng.}, vol. 47, no. 9, pp. 1185--1194, 2000.

\bibitem{physionet} A. L. Goldberger, L. A. N. Amaral, L. Glass, J. M. Hausdorff, P. Ch. Ivanov,
R. G. Mark, J. E. Mietus, G. B. Moody, C.-K. Peng, and H. E. Stanley, ``PhysioBank, PhysioToolkit,
and PhysioNet: Components of a new research resource for complex physiologic signals,''
\emph{Circulation}, vol. 101, no. 23, pp. e215--e220, 2000.

\bibitem{dod} A. Guillot, F. Sauvet, E. H. During, and V. Thorey, ``Dreem Open Datasets: Multi-scored
sleep datasets to compare human and automated sleep staging,'' \emph{IEEE Trans. Neural Syst.
Rehabil. Eng.}, vol. 28, no. 9, pp. 1955--1965, 2020.

\bibitem{causalattennet} J. Pan, Y. Feng, P. Zhao, X. Zou, A. Hou, and X. Che, ``CausalAttenNet: A
fast and long-term-temporal network for automatic sleep staging with single-channel EEG,''
\emph{IEEE Trans. Instrum. Meas.}, vol. 73, pp. 1--13, 2024.

\bibitem{bresch} E. Bresch, U. Gro\ss{}ekath\"ofer, and G. Garcia-Molina, ``Recurrent deep neural
networks for real-time sleep stage classification from single channel EEG,'' \emph{Front. Comput.
Neurosci.}, vol. 12, art. 85, 2018.

\bibitem{snoozenet} A. Mehrotra, B. Malladi, and A. Hanchate, ``SnoozeNet: Causality-driven
transformers for interpretable sleep stage classification,'' in \emph{Proc. 2024 IEEE MIT
Undergraduate Research Technology Conf. (URTC)}, 2024, pp. 1--5.

\bibitem{microsleepnet} G. Liu, G. Wei, S. Sun, D. Mao, J. Zhang, D. Zhao, X. Tian, W. Wang, and
N. Chen, ``Micro SleepNet: Efficient deep learning model for mobile terminal real-time sleep
staging,'' \emph{Front. Neurosci.}, vol. 17, art. 1218072, 2023.

\bibitem{tcn} S. Bai, J. Z. Kolter, and V. Koltun, ``An empirical evaluation of generic convolutional
and recurrent networks for sequence modeling,'' \emph{arXiv preprint arXiv:1803.01271}, 2018.

\bibitem{aasm} R. B. Berry, R. Brooks, C. E. Gamaldo, S. M. Harding, R. M. Lloyd, S. F. Quan,
M. V. Troester, and B. V. Vaughn, \emph{The AASM Manual for the Scoring of Sleep and Associated
Events: Rules, Terminology and Technical Specifications, Version 2.6}. Darien, IL: American Academy
of Sleep Medicine, 2020.

\bibitem{utime} M. Perslev, M. Jensen, S. Darkner, P. J. Jennum, and C. Igel, ``U-Time: A fully
convolutional network for time series segmentation applied to sleep staging,'' in \emph{Advances in
Neural Information Processing Systems 32 (NeurIPS 2019)}, 2019.

\bibitem{usleep} M. Perslev, S. Darkner, L. Kempfner, M. Nikolic, P. J. Jennum, and C. Igel,
``U-Sleep: Resilient high-frequency sleep staging,'' \emph{npj Digit. Med.}, vol. 4, art. 72, 2021.

\bibitem{tinysleepnet} A. Supratak and Y. Guo, ``TinySleepNet: An efficient deep learning model for
sleep stage scoring based on raw single-channel EEG,'' in \emph{Proc. 42nd Annu. Int. Conf. IEEE Eng.
Med. Biol. Soc. (EMBC)}, 2020, pp. 641--644.

\bibitem{seqsleepnet} H. Phan, F. Andreotti, N. Cooray, O. Y. Ch\'en, and M. De Vos, ``SeqSleepNet:
End-to-end hierarchical recurrent neural network for sequence-to-sequence automatic sleep staging,''
\emph{IEEE Trans. Neural Syst. Rehabil. Eng.}, vol. 27, no. 3, pp. 400--410, 2019.

\bibitem{sleepeegnet} S. Mousavi, F. Afghah, and U. R. Acharya, ``SleepEEGNet: Automated sleep stage
scoring with sequence to sequence deep learning approach,'' \emph{PLoS ONE}, vol. 14, no. 5,
art. e0216456, 2019.

\bibitem{iitnet} H. Seo, S. Back, S. Lee, D. Park, T. Kim, and K. Lee, ``Intra- and inter-epoch
temporal context network (IITNet) using sub-epoch features for automatic sleep scoring on raw
single-channel EEG,'' \emph{Biomed. Signal Process. Control}, vol. 61, art. 102037, 2020.

\bibitem{sors} A. Sors, S. Bonnet, S. Mirek, L. Vercueil, and J.-F. Payen, ``A convolutional neural
network for sleep stage scoring from raw single-channel EEG,'' \emph{Biomed. Signal Process.
Control}, vol. 42, pp. 107--114, 2018.

\bibitem{attention} A. Vaswani, N. Shazeer, N. Parmar, J. Uszkoreit, L. Jones, A. N. Gomez,
L. Kaiser, and I. Polosukhin, ``Attention is all you need,'' in \emph{Advances in Neural Information
Processing Systems 30 (NeurIPS 2017)}, 2017, pp. 5998--6008.

\bibitem{focal} T.-Y. Lin, P. Goyal, R. Girshick, K. He, and P. Doll\'ar, ``Focal loss for dense
object detection,'' in \emph{Proc. IEEE Int. Conf. Comput. Vis. (ICCV)}, 2017, pp. 2980--2988.

\bibitem{resnet} K. He, X. Zhang, S. Ren, and J. Sun, ``Deep residual learning for image
recognition,'' in \emph{Proc. IEEE Conf. Comput. Vis. Pattern Recognit. (CVPR)}, 2016, pp. 770--778.

\bibitem{batchnorm} S. Ioffe and C. Szegedy, ``Batch normalization: Accelerating deep network
training by reducing internal covariate shift,'' in \emph{Proc. 32nd Int. Conf. Mach. Learn. (ICML)},
PMLR vol. 37, 2015, pp. 448--456.

\bibitem{cambrain} \emph{CaMBRAIN: Real-time, continuous EEG inference with causal state space
models}, arXiv preprint arXiv:2605.28792, 2026.

\end{thebibliography}

\end{document}
